(1) 守恒角动量 $L=mr^2\dot\phi$，径向能量为 $E=m\dot r^2/2+L^2/(2mr^2)+U(r)$。$s\ne1$ 时 $U=kr^{1-s}/(1-s)$，而 $s=1$ 时 $U=k\ln r$，均满足 $F_r=-U'(r)=-kr^{-s}$。

(2) 圆轨道条件 $L^2/(mr_0^3)=kr_0^{-s}$ 给出 $L^2=mkr_0^{3-s}$ 与 $\Omega_0^2=k/(mr_0^{s+1})$。有效势二阶导为 $U_{\rm eff}''(r_0)=(3-s)kr_0^{-s-1}$，其正负控制径向线性稳定性。

(3) 令 $r=r_0+\delta r$，有 $m\delta\ddot r=-U_{\rm eff}''(r_0)\delta r$，从而 $\kappa^2=(3-s)\Omega_0^2$。$s<3$ 稳定，$s>3$ 指数失稳，$s=3$ 为线性边界、频率为零，必须检查精确有效势与非线性项。

(4) 近圆极限中每一径向周期 $T_r=2\pi/\kappa$ 的方位进动角 $\Delta\phi=\Omega_0T_r=2\pi/\sqrt{3-s}$。相对Kepler椭圆每次近心经过的多余进动为 $2\pi(1/\sqrt{3-s}-1)$；$s=2$ 时径向与角频率相等，恰无一阶进动。
